% Stage E: the side of the denominators: the 2-adic information with its classes modulo 64, the
% class word as the starting value of the first stage and the 2-adic exit; for odd degree the
% squares and the divisors of the leading coefficient times squares (setup_us1) with the valuation
% test of the slopes; for even degree the forbidden divisors, tested with words of 64 denominators
% and by their valuations, and the Jacobi test, by the binary algorithm and by tables; and the loop
% over the denominators in the four shapes of run_setup and run_block.
\chapter{The denominators}\label{chap:denominators}

Stage~E models what ratpoints does on the side of the denominators. Before the sieve it computes,
for each class of the denominators modulo $64$, which numerators modulo $64$ the class admits, and
which classes admit one at all; for a polynomial of odd degree it restricts the denominators to
squares, or to squares times squarefree divisors of the leading coefficient with a test on their
valuations; for even degree it finds primes that divide no denominator of a point, or only to
certain powers, and tests the denominators by a Jacobi symbol; and it runs the loop over the
denominators in one of four shapes that combine these tests. Each test only removes candidates that
are not points, so the points found do not change. What changes is their order: for a polynomial
searched of odd degree whose leading coefficient is not $\pm 1$, ratpoints visits the denominators
divisor by divisor when the factorization of the leading coefficient over its table of primes is
complete and gives fewer than $512$ squarefree divisors up to the largest denominator, so that
\verb|ratpoints '1 2' 20 -q| prints the points of the denominator~$9$ before those of the
denominator~$2$: the divisors of the leading coefficient $2$ are $1$ and $2$, and the loop takes $1,
4, 9, 16$ and then $2, 8, 18$.

So the output is the points at infinity followed by the points of the brute-force reference
over the denominators that the loop of the search visits, in its order, each mapped back after a
reversal. It is a permutation of the reference, as in stage~D, so that the main theorems, the
refusal theorems and the statements that comparator certifies do not change; and it is the reference
itself, without a reversal, wherever the loop visits the denominators in increasing order, in
particular for every shape but the divisors times squares. Besides the deviations of
Section~\ref{sec:denDeviations}, which name their stages, the model of this chapter keeps two
simplifications of the loop, each the subject of a later stage: every class of the denominators
packs its numerators with the stride~$1$, so that the class word is a pattern of the plain layout of
the model, where ratpoints packs the numerators of a class that are all congruent modulo $2^e$ into
the bits of every $2^e$-th numerator (stage~F3); and the loop visits the denominators of each
segment in order, where ratpoints cuts the segments into blocks for its threads and for the
correction of the moduli between blocks (stage~I). Neither changes the points or their order.

\section{What the C code does on the side of the denominators}\label{sec:denC}

For the polynomial searched, of degree $d$ with the coefficients $c_0, \dots, c_d$ that the
set-up of stage~D hands on, the form $F(x, z) = \sum_k c_k x^k z^{D - k}$ of even degree
$D = d + (d \bmod 2)$, and the largest denominator $b_{\mathrm{high}}$, which is $H$ until stage~H,
\verb|find_points_work_1|:
\begin{enumerate}
  \item sets the flag of the checked loop over the denominators, \verb|RATPOINTS_CHECK_DENOM|, for
    every degree (\csource{find_points.c}{4606}{4607}); clears it for odd $d$; and for even $d$ with
    $c_d$ not a square sets \verb|RATPOINTS_USE_JACOBI|, unless the caller has switched the Jacobi
    test off (\verb|RATPOINTS_NO_JACOBI|, the program's \verb|-j|)
    (\csource{find_points.c}{4900}{4907});
  \item for odd $d$ and $|c_d| = 1$, sets \verb|RATPOINTS_USE_SQUARES|: the denominator of every
    point is then a square; for odd $d$ and $|c_d| \ne 1$, calls \verb|setup_us1|
    (\csource{find_points.c}{766}{935}), which factors $|c_d|$ over $2$ and the $171$ odd primes of
    \verb|primes.h|, taking a cofactor below $p^2$, which has no prime factor below $p$, as a prime;
    when the factorization is complete, lists the squarefree divisors of $c_d$ up to
    $b_{\mathrm{high}}$ in the order of binary counting over its primes in increasing order, the
    first prime the lowest bit; and when there are fewer than $\mathrm{maxDivisors} = 512$ of them,
    computes for each prime $p$ of $c_d$ the slope of a valuation test and sets
    \verb|RATPOINTS_USE_SQUARES1| (\csource{find_points.c}{4909}{4943});
  \item computes the 2-adic information (\verb|get_2adic_info|, \csource{find_points.c}{981}{1115}):
    the coefficients modulo $64$; the residues $k$ modulo $64$ at which $f(k)$ is a square modulo
    $64$, and the even residues $t$ at which the reversed polynomial
    $\tilde f(t) = t^D f(1/t)$ is one, both by Horner's rule in the arithmetic of
    \verb|unsigned long|; from them, by scattering, for each class $b$ modulo $64$ the word
    $\mathrm{nb}[b]$ of the admissible numerators modulo $64$, bit $a$ set exactly when $a$ or $b$
    is odd and $F(a, b)$ is a square modulo $64$; the packing of each class, the largest
    $e \le 6$ for which its admissible numerators lie in one class $a_0$ modulo $2^e$, and its
    pattern in that packing; and the word $\mathrm{den\_bits}$ of the classes that admit a
    numerator;
  \item ends the search with no points when no class admits a numerator and there is no point at
    infinity (\csource{find_points.c}{4954}{4963}), which happens only for even $d$, $c_0$ and
    $c_d$ non-squares modulo $64$, and never after a reversal;
  \item in \verb|sieving_info|, for each prime $p$ that it examines and while the checked loop is on
    and fewer than \verb|max_forbidden| forbidden divisors have been found: for $p \mid c_d$, the
    mask of the valuations $m \ge 1$ with $p^m \le b_{\mathrm{high}}$ that no denominator of a point
    has, by the terms of $F$ of least valuation at $p$ (\verb|forbidden_valuations|,
    \csource{find_points.c}{2192}{2234}), recorded, when it is not zero, as a prime tested with
    words of $64$ denominators when it holds every such valuation, and as a valuation entry
    otherwise; for $p \nmid c_d$ with $c_d$ not a square modulo $p$, a prime tested with words
    (\csource{find_points.c}{2624}{2677}); then, unless $c_d$ is a square, the further primes of the
    table with $p^2 \le b_{\mathrm{high}}$ and the Kronecker symbol $(c_d / p) = -1$, tested with
    words, whose words it makes for the primes beyond the table of the build
    (\csource{find_points.c}{2709}{2789}); and it clears the flag of the checked loop when there is
    no forbidden divisor and no Jacobi test (\csource{find_points.c}{2793}{2795});
  \item chooses the shape of the loop (\verb|run_setup|, \csource{find_points.c}{3778}{3877}): the
    squares $k^2$ under \verb|RATPOINTS_USE_SQUARES|; the divisors times squares $s k^2$ under
    \verb|RATPOINTS_USE_SQUARES1|, one segment per divisor, in the order of the divisors; the words
    of $64$ denominators under \verb|RATPOINTS_CHECK_DENOM|, with the tables of the Jacobi test
    when the Jacobi test is on (\verb|RATPOINTS_USE_JACOBI|), $b_{\mathrm{high}} < 2^{32}$ and the
    odd part of $c_d$ factors over the table into at most $16$ primes whose tables of the
    non-squares, for the odd exponents, take at most $8192$ entries (\verb|jacobi_setup|,
    \csource{find_points.c}{683}{731}); and
    every denominator otherwise; and cuts the segments into blocks, which the calling thread or the
    threads sieve;
  \item in each block (\verb|run_block|, \csource{find_points.c}{3889}{4078}) visits the
    denominators of its shape in order and tests each: for the squares, the divisors and every
    denominator, that the pattern of the class of $b \bmod 64$ is not empty, and for the divisors
    then the valuation test of the slopes; for the words, it and-s the word of
    $\mathrm{den\_bits}$ with the words of the primes tested with words and takes the set bits in
    increasing order, each then tested by the Jacobi test by tables, the valuation entries, and,
    with the Jacobi test on and no tables, the Jacobi test by the binary algorithm, on \verb|long|s
    (\verb|jacobi1|) or
    on GMP integers (\verb|jacobi|); a denominator that passes is sifted, and the first stage of the
    sieve starts from the pattern of its class rather than from all ones
    (\csource{sift.c}{948}{955}).
\end{enumerate}
So the program prints the points at infinity, then the points of each denominator in the order of
the loop of its shape, the numerators of each denominator in increasing order: in increasing order
of the denominators for the squares, the words and every denominator, and divisor by divisor for the
divisors times squares. The tests and the 2-adic information only remove denominators and
numerators that carry no point; the blocks, the threads and the correction of the moduli between
the blocks change neither the points nor their order; and the packing of the classes changes only
which bit of a bit array stands for which numerator.

\section{The 2-adic information}\label{sec:twoAdic}

\begin{definition}[The 2-adic information]\label{def:twoAdic}
  \uses{def:ctz}
  \lean{Ratpoints.Sieve.squares64, Ratpoints.Sieve.coeffsMod64, Ratpoints.Sieve.fsq, Ratpoints.Sieve.gsq, Ratpoints.Sieve.twoAdic, Ratpoints.Sieve.denBits}
  \leanok
  The squares modulo $64$ as a word, bit $s$ set for the twelve squares
  $0, 1, 4, 9, 16, 17, 25, 33, 36, 41, 49, 57$ (\csource{find_points.c}{981}{981}); the
  coefficients modulo $64$ as the form of degree $D$ reads them, with $c_D = 0$ for odd $d$
  (\csource{find_points.c}{1002}{1004}); the word $\mathrm{fsq}$ of the residues $k$ with
  $f(k)$ a square modulo $64$ (\csource{find_points.c}{1006}{1014}) and the word $\mathrm{gsq}$ of
  the even residues $t$ with $\tilde f(t)$ one (\csource{find_points.c}{1015}{1021}), by Horner's
  rule in the arithmetic of $64$-bit words, which is exact modulo $64$; the words
  $\mathrm{nb}[b]$, $b < 64$: from zero, for each set bit $k$ of $\mathrm{fsq}$, lowest first, bit
  $k b \bmod 64$ of the word of every odd $b$, and for each set bit $t$ of $\mathrm{gsq}$, bit $a$
  of the word of $t a \bmod 64$ for every odd $a$ (\csource{find_points.c}{983}{994},
  \csource{find_points.c}{996}{1000}, \csource{find_points.c}{1023}{1050},
  \csource{find_points.c}{1106}{1115}); and $\mathrm{den\_bits}$, bit $b$ set exactly when
  $\mathrm{nb}[b] \ne 0$ (\csource{find_points.c}{995}{995}, \csource{find_points.c}{1064}{1065},
  \csource{find_points.c}{1070}{1070}, \csource{find_points.c}{1096}{1097},
  \csource{find_points.c}{1102}{1104}). The model keeps the stride~$1$ for every class, the
  stride exponent $0$ and the offset $a_0 = 0$, so that the pattern of a class is its word
  $\mathrm{nb}[b]$ itself, and computes $\mathrm{den\_bits}$ from the words. Bit $a$ of
  $\mathrm{nb}[b]$ is set exactly when $a$ or $b$ is odd and $F(a, b)$ is a square modulo $64$.
\end{definition}

\begin{lemma}[The class of a point]\label{lem:twoAdic_point}
  \uses{def:twoAdic, def:binaryForm, def:bruteForce, def:andRows}
  \lean{Ratpoints.Sieve.squares64_test_iff, Ratpoints.Sieve.getLsbD_fsq_coeffsMod64, Ratpoints.Sieve.getLsbD_gsq_coeffsMod64, Ratpoints.Sieve.TwoAdicBits, Ratpoints.Sieve.twoAdic_spec, Ratpoints.Sieve.getLsbD_run_twoAdic, Ratpoints.isSquare_sum_mul_pow_sub_of_sq_eq_binaryForm, Ratpoints.Sieve.classBit_twoAdic_of_sq_eq, Ratpoints.Sieve.getLsbD_denBits, Ratpoints.Sieve.getElem!_twoAdic_ne_zero}
  \leanok
  Let $a$ and $y$ be integers and $b$ a natural number with $\gcd(a, b) = 1$ and $y^2 = F(a, b)$, as
  for a normalized point $(a, y, b)$. Then bit $a \bmod 64$ of $\mathrm{nb}[b \bmod 64]$ is set, and
  so is bit $b \bmod 64$ of $\mathrm{den\_bits}$. The contract of the words is exact: bit $a$ of
  $\mathrm{nb}[b]$ is set if and only if $b$ is odd and $a = k b \bmod 64$ for a set bit $k$ of
  $\mathrm{fsq}$, or $a < 64$ is odd and $b = t a \bmod 64$ for a set bit $t$ of $\mathrm{gsq}$; and
  bit $k$ of $\mathrm{fsq}$ is set exactly when $f(k)$ is a square in $\mathbb{Z}/64$, bit $t$ of
  $\mathrm{gsq}$ exactly when $t$ is even and $\tilde f(t)$ is one.
\end{lemma}

\begin{proof}
  \uses{lem:isSquare_div}
  \leanok
  Since $\gcd(a, b) = 1$, $a$ or $b$ is odd. For odd $b$, $F(a, b) = b^D f(a b^{-1})$ modulo $64$,
  where $b^D$ is the square of a unit since $D$ is even; so $f(k)$ is a square modulo $64$ for $k =
  a b^{-1} \bmod 64$, bit $k$ of $\mathrm{fsq}$ is set, and the scatter sets bit $k b \equiv a$ of
  $\mathrm{nb}[b \bmod 64]$. For even $b$, $a$ is odd and $F(a, b) = a^D \tilde f(b a^{-1})$ modulo
  $64$; so bit $t = b a^{-1} \bmod 64$ of $\mathrm{gsq}$ is set, $t$ is even, and the scatter sets
  bit $a \bmod 64$ of $\mathrm{nb}[t a \bmod 64]$, the word of $b \bmod 64$. Horner's rule on
  $64$-bit words, read in $\mathbb{Z}/64$, is Horner's rule there; the twelve squares of the word
  $\mathrm{squares64}$ are checked by \verb|decide|.
\end{proof}

The packing by the stride, from the masks of the classes to the stride and the offset of each class
(\csource{find_points.c}{1058}{1062}, \csource{find_points.c}{1066}{1067},
\csource{find_points.c}{1069}{1069}, \csource{find_points.c}{1071}{1095},
\csource{find_points.c}{1099}{1101}), is stage~F3. The first stage of the sieve starts from the word
of the class of the denominator: the components of Chapter~\ref{chap:words} that make a piece and
the numerators of one denominator take the class word $\omega$ as an argument, and the first stage
(Definition~\ref{def:andRows}) writes $\omega$ into every word of the piece in place of all ones, as
ratpoints writes the pattern of the class, which \verb|get_2adic_info| makes once per class, one
word repeated over a bit array by the macro \verb|RBA| of the build, and copies into each bit array;
so the node of the first stage holds the line that makes the pattern and the macro of every build
(\csource{find_points.c}{1098}{1098}, \csource{rp-private.h}{129}{132},
\csource{rp-private.h}{203}{204}, \csource{rp-private.h}{235}{235},
\csource{rp-private.h}{261}{261}, \csource{rp-private.h}{285}{285}); and their contracts state the
class bit: bit $t$ of word $y$ of a piece is bit $t$ of $\omega$ and the survivor test of the
numerator $R w_0 + 64 y + t$, which is congruent to $t$ modulo $64$ since $64 \mid R$
(\verb|Ratpoints.Sieve.classBit|, the bit $a \bmod 64$ of a word, and
\verb|Ratpoints.Sieve.classBit_rbaLength_mul_add|, which the node of the first stage names). So
\verb|sift| with the class word $\omega$ gives what the brute-force reference finds over $b$ and the
numerators that pass bit $a \bmod 64$ of $\omega$ and survive the sieve entries
(Definitions~\ref{def:siftPiece} and~\ref{def:sift}): the equation of
Lemma~\ref{lem:sift_eq_checkDen} filters the points of the plain sieve by the class bit of their
numerators (\verb|Ratpoints.BruteForce.pointsOfDen_filter|, which its node names); with the word of
the class of $b$, by Lemma~\ref{lem:twoAdic_point}, that is what the reference finds over $b$ and
the numerators that survive the sieve entries.

\begin{lemma}[The 2-adic exit]\label{lem:twoAdic_exit}
  \uses{lem:twoAdic_point, def:IsPoint}
  \lean{Ratpoints.Sieve.not_isPoint_of_denBits_eq_zero}
  \leanok
  If $\mathrm{den\_bits} = 0$, then the curve has no normalized point of any height, a point at
  infinity included. So the test of the points at infinity at the exit of the C code changes
  nothing: $\mathrm{den\_bits} = 0$ happens only for even $d$ with $c_d$ not a square modulo $64$.
\end{lemma}

\begin{proof}
  \leanok
  A point $(a, y, b)$, $b = 0$ included, has bit $a \bmod 64$ of $\mathrm{nb}[b \bmod 64]$ set, and
  so bit $b \bmod 64$ of $\mathrm{den\_bits}$ (Lemma~\ref{lem:twoAdic_point}).
\end{proof}

\section{Odd degree: squares and divisors}\label{sec:oddDegree}

For odd $d$, $F(a, b) = b \, G(a, b)$ with
$G(a, b) = c_d a^d + b \sum_{k < d} c_k a^k b^{d - 1 - k}$, and for coprime $a$ and $b$,
$\gcd(b, G(a, b)) = \gcd(b, c_d)$. This restricts the denominators of the points.

\begin{definition}[The valuation]\label{def:valuation}
  \lean{Ratpoints.veryBig, Ratpoints.Setup.valuation, Ratpoints.Setup.valuation_spec, Ratpoints.Setup.fst_run_valuation, Ratpoints.Setup.snd_run_valuation}
  \leanok
  The valuation $v$ of an integer $n \ne 0$ at $p$ and the unit part $|n| / p^v$ modulo $p$, by
  repeated division with remainder, and $(\mathrm{veryBig}, 0)$ for $n = 0$, with
  $\mathrm{veryBig} = 1000$ (\verb|VERY_BIG|, \csource{find_points.c}{422}{422}): one function for
  \verb|valuation| on GMP integers, which also returns the unit part, and \verb|valuation1| on
  \verb|long|s (\csource{find_points.c}{424}{454}). For $p > 1$ and $n \ne 0$, $v$ is the $p$-adic
  valuation of $n$.
\end{definition}

\begin{definition}[The divisors of the leading coefficient]\label{def:setupUs1}
  \uses{def:valuation, def:primeTable, def:useCLong}
  \lean{Ratpoints.maxDivisors, Ratpoints.Setup.DenInfo, Ratpoints.Setup.factorLcf, Ratpoints.Setup.squarefreeDivisors, Ratpoints.Setup.slope, Ratpoints.Setup.setupUs1, Ratpoints.Setup.IsFactorization, Ratpoints.Setup.factorLcf_spec, Ratpoints.Setup.extendDivisors, Ratpoints.Setup.squarefreeDivisors_spec, Ratpoints.Setup.slope_spec, Ratpoints.Setup.setupUs1_spec, Ratpoints.Setup.isFactorization_of_run_factorLcf_eq, Ratpoints.Setup.exists_of_setupUs1_eq_some, Ratpoints.Setup.eq_foldl_of_setupUs1_eq_some, Ratpoints.mem_primeTable_of_prime, Ratpoints.mem_of_toList_primeTable_eq, List.mem_of_mem_append_cons_of_rel, Nat.prime_of_lt_mul_self, Ratpoints.Setup.fitsLong_mul_of_le_div, Ratpoints.Setup.mul_self_lt_of_mem_primeTable, Ratpoints.Setup.fitsLong_of_lt_mul_self}
  \leanok
  For odd $d$ and $|c_d| \ne 1$, as \verb|setup_us1| computes it: the factorization of $|c_d|$, the
  power of $2$ divided off, then for the odd primes $p$ of the table in turn, while the cofactor is
  not $1$, a cofactor below $p^2$ taken as a prime with the valuation $1$, and otherwise the power
  of $p$ divided off (\csource{find_points.c}{777}{860}); complete exactly when the cofactor ends at
  $1$, that is when the part of $|c_d|$ prime to $2 \cdot 3 \cdots 1019$ is below $1021^2$ or a
  power of $1021$; the squarefree divisors at most $b_{\mathrm{high}}$, from $1$, each prime in turn
  appending $x p$ for every divisor $x \le \lfloor b_{\mathrm{high}} / p \rfloor$ found so far, in
  their order, and stopping at $\mathrm{maxDivisors} = 512$ entries (\verb|MAX_DIVISORS|,
  \csource{find_points.c}{121}{121}; \csource{find_points.c}{864}{880}); and for each prime $p$ of
  $c_d$ a record of $p$, its valuation $v_p(c_d)$ and the slope (\verb|use_squares1_info|,
  \csource{find_points.c}{155}{155}), the least $n \ge 1$ with $(d - k) n > v_p(c_d) - v_p(c_k)$ for
  every $k < d$, with $v_p(0) = \mathrm{veryBig}$ (Definition~\ref{def:valuation}), so that a zero
  coefficient imposes nothing while $v_p(c_d) \le \mathrm{veryBig}$, by the macro \verb|CEIL|
  (\csource{find_points.c}{894}{912}, \csource{rp-private.h}{87}{87}); none when the factorization
  is not complete or there are $512$ divisors (\csource{find_points.c}{766}{776},
  \csource{find_points.c}{863}{863}, \csource{find_points.c}{881}{893},
  \csource{find_points.c}{914}{935}). The primes come in increasing order, so the divisors are the
  products of the subsets of the primes that are at most $b_{\mathrm{high}}$, in the order of their
  bit masks, each once. The contracts: a factorization with the cofactor $1$ lists the distinct
  primes of $|c_d|$ with their valuations, since the table holds every odd prime below $1022$ in
  increasing order, which \verb|decide| checks, so that a cofactor below $p^2$, every smaller prime
  divided out, is a prime; a list of fewer than $\mathrm{maxDivisors}$ divisors is the fold of the
  step that appends $x p$ for the divisors $x \le \lfloor b_{\mathrm{high}} / p \rfloor$; the slope
  $n$ is at least $1$ with $(d - k) n \ge 1 + v_p(c_d) - v_p(c_k)$ for every $k < d$, with the
  valuation of Definition~\ref{def:valuation}, since \verb|CEIL| rounds up; and the products $x p$
  are \verb|long|s, and the square $p^2$ of a prime of the table and a cofactor below it are
  \verb|int|s, so that the test of the C code whether the cofactor is a \verb|long| always passes
  and the record holds it.
\end{definition}

\begin{lemma}[A square sum with one term of least valuation]\label{lem:least_valuation}
  \lean{Int.even_and_exists_isSquare_add_of_sq_eq_sum}
  \leanok
  Let $p$ be a prime and $y^2 = \sum_{i \in s} f_i$ a sum of integers in which one term
  $f_j = p^v u$, $p \nmid u$, has the least valuation: $p^{v + 1}$ divides every other term. Then
  $v$ is even and $u + p t$ is a square for some integer $t$, so that $u$ is a square modulo $p$.
\end{lemma}

\begin{proof}
  \leanok
  The sum is $p^v (u + p t)$ for an integer $t$, and $p \nmid u + p t$; since $y^2$ has an even
  valuation, $v$ is even, and $y / p^{v / 2}$ is an integer whose square is $u + p t$.
\end{proof}

\begin{lemma}[The denominators of odd degree]\label{lem:odd_denominators}
  \uses{def:setupUs1, def:binaryForm, lem:squarefree_binaryFormPoly}
  \lean{Ratpoints.even_of_sq_eq_binaryForm_of_odd, Ratpoints.isSquare_of_sq_eq_binaryForm_of_odd, Ratpoints.exists_eq_mul_mul_of_sq_eq_binaryForm_of_odd, Ratpoints.Setup.mem_foldl_extendDivisors_one_iff, Ratpoints.Setup.nodup_foldl_extendDivisors, Ratpoints.Setup.mem_iff_of_setupUs1_eq_some, Ratpoints.Setup.nodup_of_setupUs1_eq_some, Ratpoints.even_add_of_sq_eq_binaryForm_of_odd, Ratpoints.Setup.DenInfo.excludes, Ratpoints.Setup.DenInfo.excludes_eq_false_of_sq_eq}
  \leanok
  Let $d$ be odd, $a$ and $y$ integers and $b \ge 1$ a natural number with $\gcd(a, b) = 1$ and
  $y^2 = F(a, b)$, as for a normalized point $(a, y, b)$.
  \begin{enumerate}
    \item If $|c_d| = 1$, then $b$ is a square.
    \item $b = s k^2$ with $s$ squarefree and every prime of $s$ dividing $c_d$; when the divisors
      of Definition~\ref{def:setupUs1} exist and $b \le b_{\mathrm{high}}$, $s$ is one of them.
    \item For every prime $p$ of $c_d$ with its slope $n$: if $v_p(b) \ge n$, then
      $v_p(b) + v_p(c_d)$ is even; so no record of the divisors excludes $b$.
  \end{enumerate}
\end{lemma}

\begin{proof}
  \uses{lem:least_valuation}
  \leanok
  Let $p \nmid c_d$ be a prime with $b = p^m b'$, $m \ge 1$ and $p \nmid b'$. Then $p \nmid a$, and
  the term $c_d a^d b$ of $F(a, b)$ alone has the least valuation $m$, every other term
  $c_k a^k b^{d + 1 - k}$ having at least $2m$; so $m$ is even (Lemma~\ref{lem:least_valuation}).
  Hence every prime of odd valuation in $b$ divides $c_d$: for $|c_d| = 1$, $b$ is a square, and in
  general $b = s k^2$ with $s$ squarefree and its primes among those of $c_d$. The factorization is
  that of $|c_d|$, since a cofactor below $p^2$ without prime factors below $p$ is a prime (the
  table holds every odd prime below $1022$ in increasing order); with fewer than $512$ divisors, the
  list holds every product of distinct primes of $c_d$ at most $b_{\mathrm{high}}$, each once, and
  $s \le b$ is one. For (3): $p \nmid a$; the term $c_d a^d$ of $G$ has the valuation $v_p(c_d)$,
  and a term $c_k a^k b^{d - k}$ with $k < d$ has a larger one,
  $v_p(c_k) + (d - k) v_p(b) \ge v_p(c_k) + (d - k) n > v_p(c_d)$; so $v_p(G(a, b)) = v_p(c_d)$, and
  $v_p(b) + v_p(c_d) = v_p(y^2)$ is even.
\end{proof}

\section{Even degree: the forbidden divisors}\label{sec:forbidden}

For even $d$ and a prime $p$, write $F(a, b) = \sum_{j \le d} c_{d - j} a^{d - j} b^j$; for a
point with $v_p(b) = m \ge 1$, $p \nmid a$, and the term $j$ has the valuation $w_j + j m$, where
$w_j = v_p(c_{d - j})$.

\begin{definition}[The forbidden valuations]\label{def:forbiddenValuations}
  \uses{def:valuation, def:squaresMod}
  \lean{Ratpoints.Sieve.unitParts, Ratpoints.Sieve.forbiddenAt, Ratpoints.Sieve.forbiddenValuations, Ratpoints.Sieve.allValuations, Ratpoints.Sieve.IsUnitPart, Ratpoints.Sieve.unitParts_spec, Ratpoints.Sieve.forbiddenAt_spec, Ratpoints.Sieve.forbiddenValuations_spec, Ratpoints.Sieve.allValuations_spec, Ratpoints.Sieve.lt_forty_of_pow_le, Ratpoints.Sieve.run_forbiddenValuations_lt, Ratpoints.Sieve.run_allValuations_lt}
  \leanok
  For a prime $p \mid c_d$ and the coefficients modulo $p$, which \verb|examine_prime| hands back to
  its caller: the valuation $w_j$ and the unit part $r_j$ modulo $p$, its sign included, of each
  nonzero $c_{d - j}$, the valuation $0$ and the residue of $c_{d - j}$ when that is not $0$
  (\csource{find_points.c}{2198}{2200}, \csource{find_points.c}{2204}{2215}); for $m = 1, 2, \dots$
  while $p^m \le b_{\mathrm{high}}$, bit $m$ of the mask, a natural number, set when the least
  $w_j + j m$ is taken at one $j_0$ only and is odd, or $j_0$ is even and $r_{j_0}$ is not a square
  modulo $p$ (\csource{find_points.c}{2218}{2228}; \csource{find_points.c}{2192}{2197},
  \csource{find_points.c}{2201}{2203}, \csource{find_points.c}{2216}{2217},
  \csource{find_points.c}{2229}{2234}); and the mask of every $m \ge 1$ with
  $p^m \le b_{\mathrm{high}}$ (\csource{find_points.c}{2632}{2641}). The contracts: the entries
  $(w_j, r_j)$ describe the coefficients at $p$ when the coefficients modulo $p$ are the residues of
  $c_0, \dots, c_d$; bit $m$ of the mask is set exactly when $m \ge 1$, $p^m \le b_{\mathrm{high}}$
  and the test at $m$ excludes it, which it does only in the two cases above; for an odd prime and a
  \verb|long| $b_{\mathrm{high}}$ every such $m$ is below $40$, so that the bound $m < 63$ of the C
  code never binds and both masks are below $2^{40}$.
\end{definition}

\begin{lemma}[The forbidden valuations are sound]\label{lem:forbiddenValuations}
  \uses{def:forbiddenValuations, def:binaryForm}
  \lean{Ratpoints.binaryForm_eq_sum_of_even, Ratpoints.even_and_isSquare_of_sq_eq_binaryForm_of_even, Ratpoints.Sieve.run_forbiddenAt_eq_false_of_sq_eq, Ratpoints.Sieve.testBit_run_forbiddenValuations_eq_false_of_sq_eq}
  \leanok
  For even $d$, a prime $p$, integers $a$ and $y$ and a natural number $b$ with $\gcd(a, b) = 1$ and
  $y^2 = F(a, b)$, as for a normalized point $(a, y, b)$: bit $v_p(b)$ of the mask of forbidden
  valuations at $p$, from the residues of the coefficients modulo $p$, is clear.
\end{lemma}

\begin{proof}
  \uses{lem:least_valuation}
  \leanok
  Let $m = v_p(b)$; bit $0$ is never set. If the least valuation $\mu$ of the terms is taken at
  $j_0$ only, then $v_p(F(a, b)) = \mu$ and $F(a, b) / p^{\mu}$ is congruent to
  $r_{j_0} a^{d - j_0} (b / p^m)^{j_0}$ modulo $p$; since $F(a, b) = y^2$, $\mu$ is even and this
  unit is a square (Lemma~\ref{lem:least_valuation}), and for even $j_0$,
  $a^{d - j_0} (b / p^m)^{j_0}$ is the square of a unit, so $r_{j_0}$ is a square modulo $p$. The
  argument is that of the Newton polygon at $p$; it does not need $p \mid c_d$, the case in which
  the C code computes the mask.
\end{proof}

\begin{definition}[The forbidden divisors]\label{def:forbidden}
  \uses{def:forbiddenValuations, def:squaresMod, def:primeTable, def:reversal, def:examinePrime}
  \lean{Ratpoints.Sieve.ForbiddenEntry, Ratpoints.Sieve.ForbiddenVal, Ratpoints.Sieve.forbiddenRow, Ratpoints.Sieve.forbiddenOfPrime, Ratpoints.Sieve.forbidden, Ratpoints.Sieve.getLsbD_getElem!_forbiddenRow, Ratpoints.Sieve.IsForbiddenEntry, Ratpoints.Sieve.IsForbiddenVal, Ratpoints.Sieve.forbidden_spec}
  \leanok
  A prime tested with words, with its $p$ words, bit $j$ of word $r$ clear exactly when
  $p \mid 64 r + j$ (\csource{find_points.c}{157}{161}; the words as the build generates them for
  the primes of its table and as \verb|sieving_info| makes them beyond it,
  \csource{gen_find_points_h.c}{102}{134}, \csource{find_points.c}{2773}{2777}), and a prime with
  the mask of its excluded valuations (\csource{find_points.c}{166}{166}); for even $d$, the
  forbidden divisor of a prime $p$ from the coefficients modulo $p$ and the entry at infinity that
  its table hands back: for $p \mid c_d$, the coefficient $c_d$ being $0$ modulo $p$, with a nonzero
  mask, $p$ with its words when the mask holds every valuation that a denominator can have, and $p$
  with its mask otherwise; for $p \nmid c_d$, $p$ with its words when there is no point at infinity
  modulo $p$, that is when $c_d$ is not a square modulo $p$ (\csource{find_points.c}{2627}{2631},
  \csource{find_points.c}{2642}{2676}), which the loop of the tables records for each examined
  prime while fewer than $\mathrm{maxForbidden}$ have been found
  (Definition~\ref{def:examinePrimes}); then, from the lists of that loop and unless $c_d$ is a
  square, the primes of the table from the position that is the number of the examined primes on,
  while fewer than $\mathrm{maxForbidden}$ have been found and $p^2 \le b_{\mathrm{high}}$, each
  with its words when $(c_d / p) = -1$ (\csource{find_points.c}{2533}{2536},
  \csource{find_points.c}{2561}{2562}, \csource{find_points.c}{2709}{2772},
  \csource{find_points.c}{2778}{2789}). The contract: each entry tested with words that the two
  loops may record is the forbidden divisor of a prime of the table, from the residues of the
  coefficients and the entry at infinity of its reference table, or a prime of the table at which
  $c_d$ is not a square, with its words; each valuation entry is the forbidden divisor of a prime
  of the table; from lists of such entries, the further primes add only such entries. It holds
  whatever the cap, the order of the entries and the primes of the table examined.
\end{definition}

\begin{lemma}[The forbidden divisors are sound]\label{lem:forbidden}
  \uses{def:forbidden, lem:forbiddenValuations, def:binaryForm, lem:examinePrime_none}
  \lean{Ratpoints.Sieve.ForbiddenVal.excludes, Ratpoints.Sieve.IsForbiddenEntry.prime, Ratpoints.Sieve.IsForbiddenEntry.words_eq, Ratpoints.Sieve.IsForbiddenEntry.not_dvd, Ratpoints.Sieve.IsForbiddenVal.excludes_eq_false}
  \leanok
  For even $d$, integers $a$ and $y$ and a natural number $b \ge 1$ with $\gcd(a, b) = 1$ and
  $y^2 = F(a, b)$, as for a normalized point $(a, y, b)$: each entry tested with words that the
  contract of Definition~\ref{def:forbidden} admits is a prime $p$ with the words of $p$, and
  $p \nmid b$ when $b \le b_{\mathrm{high}}$; and for no valuation entry that it admits is bit
  $v_p(b)$ of its mask set.
\end{lemma}

\begin{proof}
  \leanok
  For a prime $p \mid c_d$ this is Lemma~\ref{lem:forbiddenValuations}: a prime tested with words
  has every valuation $m \ge 1$ with $p^m \le b_{\mathrm{high}}$ in its mask, which covers every $b$
  with $p \mid b$. For $p \nmid c_d$ with no point at infinity modulo $p$, that is with $c_d$ not a
  square modulo $p$: if $p \mid b$, then $p \nmid a$ and $F(a, b) \equiv c_d a^d$ modulo $p$, a unit
  that is not a square, which $y^2$ is not; that $c_d$ is a square modulo every prime of $b$ that
  does not divide it is the argument of Lemma~\ref{lem:examinePrime_none}.
\end{proof}

\section{The Jacobi test}\label{sec:jacobi}

For even $d$, an odd prime $q \mid b$ with $q \nmid c_d$ makes $c_d$ a square modulo $q$, by the
argument of Lemma~\ref{lem:forbidden}; so for $b^*$, the largest divisor of $b$ prime to $2 c_d$,
the Jacobi symbol $J(c_d \mid b^*)$, the product of the Legendre symbols of its prime factors, is
$1$. The forbidden divisors test the examined primes and, beyond them, the primes of the table up
to $\sqrt{b_{\mathrm{high}}}$, at most $\mathrm{maxForbidden}$ of them in all; the Jacobi test also
rejects a denominator with an odd number of prime factors $q$, prime to $2 c_d$, with
$(c_d / q) = -1$, which these do not reach.

\begin{definition}[The Jacobi symbol by the binary algorithm]\label{def:jacobi}
  \uses{def:useCLong}
  \lean{Ratpoints.Sieve.twoSign, Ratpoints.Sieve.jacobiReduce, Ratpoints.Sieve.jacobiRound, Ratpoints.Sieve.jacobi, Ratpoints.Sieve.twoSign_eq, Ratpoints.Sieve.IsJacobiReduction, Ratpoints.Sieve.jacobiReduce_spec, Ratpoints.Sieve.jacobiRound_spec, Nat.coprimePart, Ratpoints.Sieve.jacobi_spec, Ratpoints.fitsLong_neg_ite_div_four}
  \leanok
  For $b \ge 1$ and an integer $\ell$: the factors $2$ taken out of $b$, and $u = \ell \bmod b$;
  while $u \ne 0$, a round of the binary algorithm on $n = u$ and $m = b$, which takes the factors
  $2$ out of $n$ with the sign of $(2 / m)$, and then exchanges the two, with the sign of the
  reciprocity when both are $3$ modulo $4$, and reduces the larger by the smaller, halving it until
  it is odd, in turn, until they are equal; the common value is $\gcd(u, b)$, and when it is $1$ the
  sign is the symbol, and otherwise it is divided out of $b$ and $u$ reduced modulo the new $b$ when
  it is not smaller; when $u$ is $0$ the symbol is $1$ (\csource{find_points.c}{498}{566},
  \csource{find_points.c}{568}{641}). One function for \verb|jacobi| on GMP integers and
  \verb|jacobi1| on \verb|long|s, which takes the sign of $\ell$ out first, writes
  $\ell = $ \verb|LONG_MIN| as a quarter of it, reduces $|\ell|$ modulo $b$ only for $b < 2^{59}$
  and does not reduce it after a division of $b$: other steps, the same symbol; the negation of a
  \verb|long| $\ell$, with \verb|LONG_MIN| replaced by its quarter, is a \verb|long|.
\end{definition}

\begin{lemma}[The binary algorithm computes the Jacobi symbol]\label{lem:jacobi}
  \uses{def:jacobi}
  \lean{Ratpoints.Sieve.run_jacobi, Nat.coprimePart_dvd, Nat.coprime_coprimePart, Nat.dvd_coprimePart_of_dvd}
  \leanok
  For $b \ge 1$ and every integer $\ell$, the result is $J(\ell \mid b^*)$ for the largest divisor
  $b^*$ of $b$ prime to $2 \ell$, the product of the prime factors of $b$, with multiplicity, that
  do not divide $2 \ell$.
\end{lemma}

\begin{proof}
  \leanok
  The steps keep $s \cdot J(n \mid m)$: taking out a factor $2$ of $n$ multiplies the symbol by
  $(2 / m)$; the exchange by the reciprocity of two odd numbers; the subtraction changes nothing,
  $J(m - n \mid n) = J(m \mid n)$. A round ends at $m = n = \gcd(u, b)$, and the symbol is
  $J(u \mid b)$ when that is $1$; dividing $b$ by the gcd and reducing $u$ keeps $J(\ell \mid b^*)$,
  since $b^*$ does not change. Mathlib has the symbol with its reciprocity and its value at $2$
  (\verb|jacobiSym|). The loops end by the measures $x$ for the subtraction, $n + m$ for a round
  and $b$ for the divisions; the early return of the main loop is part of its invariant.
\end{proof}

\begin{definition}[The Jacobi test by tables]\label{def:jacobiTables}
  \uses{def:valuation, def:primeTable, def:ctz, def:mulMod, lem:mulMod}
  \lean{Ratpoints.jacobiPrimes, Ratpoints.jacobiTable, Ratpoints.mulModLimit, Ratpoints.Sieve.JacobiInfo, Ratpoints.Sieve.nonSquares, Ratpoints.Sieve.jacobiSigns, Ratpoints.Sieve.jacobiSetup, Ratpoints.Sieve.jacobiTest, Ratpoints.Sieve.nonSquares_spec, Ratpoints.Sieve.legendreSym_eq_ite_getElem!_nonSquares, Ratpoints.Sieve.ite_getElem!_jacobiSigns, Ratpoints.Sieve.IsJacobiInfo, Ratpoints.Sieve.jacobiSetup_spec, Ratpoints.Sieve.jacobiTest_spec, ZMod.isSquare_natCast_iff_exists_mul_self_mod_eq, Ratpoints.mulModLimit_lt_size, Ratpoints.mulModLimit_lt_two_pow_32, Ratpoints.mul_lt_two_pow_64_of_le_mulModLimit, Ratpoints.Sieve.le_mulModLimit_of_jacobiSetup_eq_some}
  \leanok
  For $\ell = \pm 2^v \prod q_i^{e_i}$ and $b_{\mathrm{high}} \le \mathrm{mulModLimit} = 2^{32} - 1$
  (\verb|RP_MULMOD_LIMIT|, \csource{rp-private.h}{421}{421}): a record
  (\csource{find_points.c}{672}{681}) of the odd primes $q_i$ of the table, found by trial division,
  at most $\mathrm{jacobiPrimes} = 16$ of them, each with its reciprocal
  (Definition~\ref{def:mulMod}) and with the table of the non-squares modulo $q_i$ when $e_i$ is odd
  (\csource{find_points.c}{709}{713}), at most $\mathrm{jacobiTable} = 8192$ entries in all
  (\verb|RP_JACOBI_PRIMES|, \verb|RP_JACOBI_TABLE|, \csource{find_points.c}{669}{669},
  \csource{find_points.c}{670}{670}), and of eight signs, the sign for $b^* \bmod 8$ the product of
  $(-1 / b^*)$ for negative $\ell$, $(2 / b^*)^v$ and the signs of the reciprocity of the
  $q_i \equiv 3 \pmod 4$ with odd $e_i$ (\csource{find_points.c}{721}{729}); none when a prime of
  $\ell$ is beyond the table, or a bound is passed (\csource{find_points.c}{683}{708},
  \csource{find_points.c}{714}{720}, \csource{find_points.c}{730}{731}). The test of $b$: its odd
  part with every $q_i$ divided out, the sign of the result modulo $8$ combined with the table entry
  of the result modulo each $q_i$ with odd $e_i$, passing when the product is $1$; the quotient by
  $q_i$ is computed with the reciprocal that the record holds, and the remainder from the quotient
  (\csource{find_points.c}{734}{759}). The contracts describe the loops: the table of $q$ marks the
  residues $x^2 \bmod q$, $1 \le x \le q / 2$, made by the steps $2x + 1$ with one subtraction, so
  that it gives the Legendre symbol of a residue other than $0$; the tables exist only for
  $b_{\mathrm{high}} \le \mathrm{mulModLimit}$ and $\ell \ne 0$, and then describe $\ell$ as above,
  their primes are primes of the table and they hold the reciprocals of their primes; and for primes
  $q$ of the tables with $b q < 2^{64}$, where the quotient with the reciprocal is exact
  (Lemma~\ref{lem:mulMod}), the test negates the exclusive or of the sign and the entries over the
  part of $b$ free of $2$ and of the primes of the tables, the odd part taken by the lowest set bit.
  A denominator of the test is at most $\mathrm{mulModLimit}$, below $2^{32}$, and below $2^{64}$,
  so that it is a word; a value at most $\mathrm{mulModLimit}$ times a modulus $p \le 2^{32}$ is
  below $2^{64}$, so that the reductions of the test by the primes of the table, below $2^{10}$, are
  exact.
\end{definition}

\begin{lemma}[The Jacobi test]\label{lem:jacobiTest}
  \uses{def:jacobiTables, def:jacobi, lem:jacobi, def:binaryForm, lem:examinePrime_none}
  \lean{Ratpoints.Sieve.jacobiSym_eq_ite_foldl, Ratpoints.Sieve.IsJacobiInfo.prod_filter_primeFactorsList_eq, Ratpoints.Sieve.run_jacobiTest, Ratpoints.Sieve.jacobiSym_eq_one_of_sq_eq_binaryForm, Ratpoints.Sieve.jacobiSym_coprimePart_eq_one_of_sq_eq, Ratpoints.Sieve.run_jacobi_eq_one_of_sq_eq, Ratpoints.Sieve.run_jacobiTest_of_sq_eq, jacobiSym.natCast_eq_ite_mul_legendreSym}
  \leanok
  When the tables exist and $1 \le b \le b_{\mathrm{high}}$, the test of $b$ passes exactly when
  $J(\ell \mid b^*) = 1$. For even $d$, integers $a$ and $y$ and a natural number $b$ with
  $\gcd(a, b) = 1$ and $y^2 = F(a, b)$, as for a normalized point $(a, y, b)$: $J(c_d \mid b') = 1$
  for every divisor $b'$ of $b$ prime to $c_d$; so for $b \ge 1$, $J(c_d \mid b^*) = 1$, and both
  the binary algorithm and, for $b \le b_{\mathrm{high}}$, the test by tables let $b$ through.
\end{lemma}

\begin{proof}
  \leanok
  The result of the divisions is $b^*$; $J(\ell \mid b^*)$ is the product of $J(-1 \mid b^*)$ for
  negative $\ell$, $J(2 \mid b^*)^v$ and the $J(q_i \mid b^*)^{e_i}$, of which those with even
  $e_i$ are $1$ and the others, by reciprocity, the Legendre symbol of $b^* \bmod q_i$ times a sign
  that depends on $b^* \bmod 4$. The second statement is the argument before
  Definition~\ref{def:jacobi}: $c_d$ is a square modulo every prime of $b$ that does not divide it
  (Lemma~\ref{lem:examinePrime_none}); the binary algorithm computes $J(c_d \mid b^*)$
  (Lemma~\ref{lem:jacobi}).
\end{proof}

\section{The loop over the denominators}\label{sec:shapes}

\begin{definition}[The shapes of the loop]\label{def:Shape}
  \uses{def:setupUs1, def:forbidden, def:jacobi, def:jacobiTables, def:valuation, def:useCLong,
    lem:odd_denominators, lem:forbidden}
  \lean{Ratpoints.Sieve.Shape, Ratpoints.Sieve.runSetup, Ratpoints.Sieve.ceilSqrt, Ratpoints.Sieve.floorSqrt, Ratpoints.Sieve.squaresDens, Ratpoints.Sieve.squares1Dens, Ratpoints.Sieve.denInfoTest, Ratpoints.Sieve.wordBits, Ratpoints.Sieve.wordsDens, Ratpoints.Sieve.wordsTest, Ratpoints.Sieve.Shape.dens, Ratpoints.Sieve.Shape.keep, Ratpoints.Sieve.ceilSqrt_le_iff, Ratpoints.Sieve.le_floorSqrt_iff, Ratpoints.Sieve.getLsbD_wordBits, Ratpoints.Sieve.wordsDens_spec, Ratpoints.Sieve.run_denInfoTest_eq_true, Ratpoints.Sieve.run_wordsTest_eq_true, Ratpoints.Sieve.fitsLong_of_mem_squaresDens, Ratpoints.Sieve.fitsLong_of_mem_squares1Dens, Ratpoints.Sieve.fitsLong_of_mem_wordsDens}
  \leanok
  The four shapes, with their data (\csource{find_points.c}{3725}{3728},
  \csource{find_points.c}{3736}{3763}), as \verb|run_setup| chooses them from the flags of the
  search (\csource{find_points.c}{3778}{3806}, \csource{find_points.c}{3876}{3877}), each with the
  list of its denominators from $b_{\mathrm{low}}$ to $b_{\mathrm{high}}$ in the order of its loop
  and its test of a single denominator:
  \begin{itemize}
    \item the squares: $k^2$ for $k$ from $\lceil \sqrt{b_{\mathrm{low}}} \rceil$ to
      $\lfloor \sqrt{b_{\mathrm{high}}} \rfloor$ (\csource{find_points.c}{464}{471},
      \csource{find_points.c}{3766}{3774}, \csource{find_points.c}{3816}{3818},
      \csource{find_points.c}{3919}{3922}), with the class test;
    \item the divisors times squares: for each divisor $s$ in turn, $s k^2$ for the $k$ with
      $b_{\mathrm{low}} \le s k^2 \le b_{\mathrm{high}}$ (\csource{find_points.c}{3819}{3825},
      \csource{find_points.c}{3937}{3940}), with the class test and then the valuation test: $b$
      fails when for some prime $p$ of $c_d$ with its slope $n$, $v_p(b) \ge n$ and
      $v_p(b) + v_p(c_d)$ is odd (\csource{find_points.c}{3944}{3953});
    \item the words: for each word $w$ from $\lfloor b_{\mathrm{low}} / 64 \rfloor$ to
      $\lfloor b_{\mathrm{high}} / 64 \rfloor$, $\mathrm{den\_bits}$ and-ed with word $w \bmod p$ of
      each prime tested with words, the bits below $b_{\mathrm{low}}$ and above $b_{\mathrm{high}}$
      cleared in the first and the last word (\csource{find_points.c}{3981}{3988},
      \csource{find_points.c}{3990}{3990}, \csource{find_points.c}{3993}{4006}), and its set bits
      $j$, lowest first, giving $64 w + j$ (\csource{find_points.c}{3826}{3828},
      \csource{find_points.c}{3967}{3980}, \csource{find_points.c}{3989}{3989},
      \csource{find_points.c}{3991}{3991}, \csource{find_points.c}{4007}{4008},
      \csource{find_points.c}{4012}{4012}, \csource{find_points.c}{4053}{4058}), each tested by the
      Jacobi test by tables, the valuation entries, and, with the Jacobi switch on and no tables,
      the binary algorithm (\csource{find_points.c}{4010}{4011},
      \csource{find_points.c}{4013}{4041});
    \item every denominator from $b_{\mathrm{low}}$ to $b_{\mathrm{high}}$, with the class test
      (\csource{find_points.c}{3812}{3815}, \csource{find_points.c}{3829}{3837},
      \csource{find_points.c}{4060}{4062}, \csource{find_points.c}{4073}{4074}).
  \end{itemize}
  The class test is that the word $\mathrm{nb}[b \bmod 64]$ is not zero
  (\csource{find_points.c}{3923}{3925}, \csource{find_points.c}{3941}{3943},
  \csource{find_points.c}{4063}{4065}); for the words it is made by $\mathrm{den\_bits}$. The model
  lists the denominators of a shape and then tests and sifts them one after the other, where
  \verb|run_block| tests and sifts each as its loop reaches it, and has no blocks. The words of a
  forbidden divisor, the word of $64$ denominators and the walk over its set bits are exact, so
  that the list of the words is the range filtered by the class bit and the primes tested with
  words; the tests of a denominator pass it when no entry of the valuation test excludes it, the
  Jacobi tests pass it and no valuation entry excludes it. The bounds of the square shapes are the
  exact integer square roots, and the denominators of the shapes are at most $b_{\mathrm{high}}$, so
  \verb|long|s for a \verb|long| $b_{\mathrm{high}}$.
\end{definition}

\begin{lemma}[The lists of the shapes]\label{lem:shape_dens}
  \uses{def:Shape, lem:twoAdic_point, lem:odd_denominators, lem:forbidden, lem:jacobiTest}
  \lean{Ratpoints.Sieve.squaresDens_eq_filter, Ratpoints.Sieve.mem_squaresDens_iff, Ratpoints.Sieve.mem_range'_ceilSqrt_iff, Ratpoints.Sieve.mem_squares1Dens_iff, Ratpoints.Sieve.nodup_squares1Dens, Nat.eq_of_mul_sq_eq_mul_sq, Ratpoints.Sieve.toList_run_wordsDens, Ratpoints.Sieve.run_denInfoTest_of_sq_eq, Ratpoints.Sieve.getLsbD_denBits_and_all_of_sq_eq, Ratpoints.Sieve.run_wordsTest_of_sq_eq, Ratpoints.Sieve.Shape.le_and_le_of_mem_dens}
  \leanok
  Let $1 \le b_{\mathrm{low}}$. The denominators of every shape lie between $b_{\mathrm{low}}$ and
  $b_{\mathrm{high}}$. The list of the squares is that of the squares from $b_{\mathrm{low}}$ to
  $b_{\mathrm{high}}$, in increasing order. The list of the divisors times squares, for squarefree
  divisors without duplicates, has no duplicates and holds exactly the $s k^2$ from
  $b_{\mathrm{low}}$ to $b_{\mathrm{high}}$ with $s$ a divisor. The list of the words, for primes
  tested with words with their words, is that of the denominators from $b_{\mathrm{low}}$ to
  $b_{\mathrm{high}}$ whose bit in $\mathrm{den\_bits}$ is set and which no such prime divides, in
  increasing order. For integers $a$ and $y$ and a natural number $1 \le b \le H$ with
  $\gcd(a, b) = 1$ and $y^2 = F(a, b)$, as for a normalized point $(a, y, b)$ of height at most $H$,
  $b$ passes the tests of the shapes: the class test, the valuation test with the records of the
  divisors of Definition~\ref{def:setupUs1}, and for forbidden divisors that the contract of
  Definition~\ref{def:forbidden} admits with $b_{\mathrm{high}} = H$ the bit of its class in
  $\mathrm{den\_bits}$, the primes tested with words and the tests of the words.
\end{lemma}

\begin{proof}
  \leanok
  The square shape and the words are the denominators of the range that pass a test, in increasing
  order, by the bounds of the square roots and by the walk over the set bits of the words; the
  divisors times squares have no duplicates since $s k^2$ determines the squarefree $s$ and $k$.
  The class bit and the primes tested with words miss no denominator of a point
  (Lemmas~\ref{lem:twoAdic_point} and~\ref{lem:forbidden}), and the tests pass it
  (Lemmas~\ref{lem:odd_denominators}, \ref{lem:forbidden} and~\ref{lem:jacobiTest}).
\end{proof}

The run (Definition~\ref{def:sieveRun}) takes the shape, the words of the 2-adic information and the
range $b_{\mathrm{low}}, \dots, b_{\mathrm{high}}$, and lives in the model file of this chapter: for
each denominator of the list of the shape that passes its test, the points that \verb|sift| finds
from the word of its class, appended in turn (\csource{find_points.c}{3889}{3904},
\csource{find_points.c}{3917}{3918}, \csource{find_points.c}{3926}{3936},
\csource{find_points.c}{3954}{3966}, \csource{find_points.c}{4009}{4009},
\csource{find_points.c}{4042}{4052}, \csource{find_points.c}{4059}{4059},
\csource{find_points.c}{4066}{4072}, \csource{find_points.c}{4075}{4078},
\csource{find_points.c}{5120}{5120}, \csource{find_points.c}{5152}{5160}). Its contract, for an
array size $A \ge 1$ and sieve entries of moduli $p > 0$, is the concatenation, over the
denominators of the list that pass the test, of what the brute-force reference finds over the
numerators that pass the class word and survive the sieve entries, mapped by $\mathrm{swapPoint}$
when the flag of the reversal is set (\verb|Ratpoints.Sieve.sieveRun_spec|,
\verb|Ratpoints.Sieve.toList_run_sieveRun|). For sound sieve entries
(Definition~\ref{def:SieveEntry}) and the shapes of Lemma~\ref{lem:shape_dens}, the filters change
nothing, and the result is the concatenation of what the reference finds over each denominator of
the list: given a list $l$ such that the list of the shape is $l$ filtered by a test, every
denominator of $l$ that carries a point passes that test and the test of the shape, and every
numerator that gives a point passes the class word and survives the sieve entries, the run gives the
reference over $l$ (Theorem~\ref{thm:run_reference},
\verb|Ratpoints.Sieve.toList_run_sieveRun_eq_flatMap|); the points at infinity come from the search.

\section{The search and the theorems}\label{sec:denSearch}

The search of the polynomial handed on (Definition~\ref{def:searchPoints}) does what
Section~\ref{sec:denC} describes, in its order: the flags of the tests on the denominators, the
divisors of the leading coefficient for odd degree, the 2-adic information and its exit, what
\verb|sieving_info| computes (Definition~\ref{def:sievingInfo}), the tables of the examined primes
with their exit, the forbidden divisors for even degree and the sieve entries of the moduli, the
switch of the checked loop, the shape, and the points at infinity followed by the run over the
denominators from $1$ to $H$. For well-formed parameters, $H \ge 1$ and $d$ the degree of $c$, its
result is the list of the points at infinity followed by the points of the reference over the
denominators of the search (Definition~\ref{def:searchDens}), in their order, each mapped by
$\mathrm{swapPoint}$ when the flag of the reversal is set, also at both exits; its node names
\verb|Ratpoints.run_searchPoints| with this statement, uses Definitions~\ref{def:Shape}
and~\ref{def:searchDens}, Lemmas~\ref{lem:twoAdic_exit} and~\ref{lem:shape_dens} and the reduction
of the run, and holds the C ranges of the flags and of the exits, the call of \verb|sieving_info|
among them, with the frame of the block that sets up the run (\csource{find_points.c}{2793}{2795},
\csource{find_points.c}{4573}{4573}, \csource{find_points.c}{4588}{4589},
\csource{find_points.c}{4595}{4596}, \csource{find_points.c}{4603}{4604},
\csource{find_points.c}{4607}{4607}, \csource{find_points.c}{4900}{4943},
\csource{find_points.c}{4948}{4948}, \csource{find_points.c}{4954}{4963},
\csource{find_points.c}{4992}{5012}, \csource{find_points.c}{5116}{5116},
\csource{find_points.c}{5137}{5139}, \csource{find_points.c}{5167}{5167}).

\begin{definition}[The denominators of the search]\label{def:searchDens}
  \uses{def:setupUs1, def:Shape}
  \lean{Ratpoints.searchDens, Ratpoints.searchDens_eq_range'}
  \leanok
  The denominators from $1$ to $H$ in the order in which the search visits those that can carry a
  point: for odd $d$ with $|c_d| \ne 1$ and divisors of the leading coefficient
  (Definition~\ref{def:setupUs1}), the divisors times squares, divisor by divisor; and otherwise
  $1, \dots, H$, since the other shapes visit increasing lists of these that contain every
  denominator of a point.
\end{definition}

\begin{theorem}[The output of the search]\label{thm:findPoints_order}
  \uses{def:findPoints, def:searchPoints, def:searchDens, lem:shape_dens, lem:twoAdic_exit,
    lem:examinePrime_none, lem:reversal_swap}
  \lean{Ratpoints.exists_eq_of_findPoints_eq_ok, Ratpoints.perm_flatMap_squares1Dens, Ratpoints.perm_bruteForce_append_flatMap_searchDens, Ratpoints.append_flatMap_searchDens_eq_bruteForce, Ratpoints.eq_bruteForce_of_findPoints_eq_ok_of_reverse_eq_false, List.perm_flatMap_of_nodup}
  \leanok
  For well-formed parameters and any options, if $\mathrm{findPoints}(c, H, P, o)$ returns the list
  $L$, and the reversal check hands on the polynomial $c'$ of degree $d'$ with the flag of the
  reversal: $L$ is the list of the points at infinity of $c'$ followed by the points that the
  brute-force reference of $c'$ finds over the denominators of the search of $c'$
  (Definition~\ref{def:searchDens}), in their order, each mapped by $\mathrm{swapPoint}$ when the
  flag is set. The points at infinity of $c'$ and the points of its reference over the
  denominators of its search are a permutation of $\mathrm{bruteForce}(c', H)$, so that $L$ is a
  permutation of $\mathrm{bruteForce}(c, H)$ (Theorem~\ref{thm:findPoints_perm}), and $L$ is
  $\mathrm{bruteForce}(c, H)$ when the options forbid the reversal and the search of $c$ does not
  take the divisors times squares.
\end{theorem}

\begin{proof}
  \uses{thm:refusals}
  \leanok
  $L$ is the list of the search of $c'$ (Definition~\ref{def:searchPoints}), which takes an exit
  only when the curve has no points (Lemmas~\ref{lem:twoAdic_exit}
  and~\ref{lem:examinePrime_none}), and otherwise gives the points of the reference over the list of
  its shape (Lemma~\ref{lem:shape_dens}), which are those over the denominators of the search,
  since a denominator that the one list has and the other lacks carries no point. The denominators
  of the search have no duplicates, lie in $1, \dots, H$ and hold every denominator of a point
  (Lemma~\ref{lem:shape_dens}; for the shapes other than the divisors they are $1, \dots, H$), so
  the points over them are a permutation of those over $1, \dots, H$. The degree that the reversal
  check hands on is that of $c'$, and without a reversal $c' = c$ (Lemma~\ref{lem:reversal_swap}).
  When $\mathrm{findPoints}$ returns a list, its height bound is below $2^{63}$ by its refusals
  (Theorem~\ref{thm:refusals}), as the search needs.
\end{proof}

Theorem~\ref{thm:findPoints_perm} states the permutation, and Theorem~\ref{thm:findPoints}
membership, no duplicates and the set of elements, for every option. The refusals and the statements
that comparator certifies are those of stage~D in their text; the structure \verb|Params| and its
well-formedness, which stage~E left as they were, gain the number of primes examined and the
composite moduli with stage~F1, which changes what those statements quantify over
(Chapter~\ref{chap:tables}).

\section{The command-line program and the differential tests}\label{sec:denCli}

The program is that of stage~D: the options that stage~E reads were read before, \verb|-j| into
the Jacobi switch of the parameters and \verb|-F| into $\mathrm{maxForbidden}$, and the order of the
printed points is that of Theorem~\ref{thm:findPoints_order}. The differential tests compare every
point list in its order, the divisors of the leading coefficient included; in the suites that run
the C program, a run on which it prints nothing (with \verb|-q|, as at the 2-adic exit) is compared
even above the largest height, which the harness skips only when the C program prints something,
so that \verb|test3.sh| compares 16 of its 21 runs. \verb|ratpoints-lean| and the C program agree
on all 3219 comparisons of the eight default suites, 261 of the runs of \verb|test4.sh| by their
points, and on the 13 of \verb|test4high|.

\section{Deviations}\label{sec:denDeviations}

Each component of this chapter differs from the C code in a stated way, recorded in its
\verb|Deviation:| paragraph:
\begin{itemize}
  \item the stride~$1$ for every class of the denominators, the pattern of a class being its word
    repeated, for the packing of the numerators of a class by the stride $2^e$ with its offsets
        (stage~F3); the word $\mathrm{den\_bits}$ computed from the words of the classes;
  \item one function for the valuations on GMP integers and on \verb|long|s, and one for the
    binary algorithm of \verb|jacobi| and \verb|jacobi1|, which take other steps to the same
    symbol; the valuations of the factorization and of the tables of the Jacobi test by this
    function, for the loops of divisions of the C code; the bounds of the square shapes by the exact
    integer square root, for the square root in \verb|double| corrected by a step either way, with
    the same values for the arguments $n \ge 1$ that the C code passes;
  \item the examined primes of the forbidden divisors are the first primes of the table up to the
    number that the parameters examine, and the further primes start at that position, for the final
    \verb|pn_lim| of \verb|sieving_info|, which it raises when the number of primes was left to the
    library and the choice of the moduli needs more (stage~I);
  \item the Kronecker symbol $(c_d / p) = -1$ by the table of the squares modulo $p$, which marks
    $0$ as a square; the square tests of the leading coefficient, for the Jacobi switch and for the
    point at infinity of the 2-adic exit, made on the polynomial searched, where the C code reads
    the flag that it computes before the reversal check and clears in its second case: the same
    value wherever it is read;
  \item the masks of the forbidden valuations as natural numbers, without the bound $m < 63$ of the
    C code, which never binds for a \verb|long| $b_{\mathrm{high}}$ and an odd prime ($3^{40} >
    2^{63}$), so that the model stays sound for every height;
  \item lists for the arrays with their terminating zeros, natural numbers for the \verb|long|s and
    \verb|int|s, divisions for the multiplications by reciprocals (\verb|RP_MULMOD|) of the first
    word of a block modulo the forbidden divisors tested with words, which come with the blocks
    (stage~I), fresh arrays for the buffers of the run and of the arguments; the values
    that could leave their C types are proved to stay in them for a \verb|long| $b_{\mathrm{high}}$,
    by lemmas that the nodes of their components name: the denominators $k^2$, $s k^2$ and
    $64 w + j$ of the shapes, the products $x p$ of the divisors and $p^m$ of the masks, the square
    $p^2$ of a prime of the table and the cofactor of the factorization below it, \verb|int|s, the
    masks, below $2^{40}$, and the negation of the quarter of \verb|LONG_MIN| in \verb|jacobi1|, the
    bounds of $p^m$ and $p^2$ in the loops of the forbidden divisors being those of the node of the
    divisors (Definition~\ref{def:setupUs1});
    and a denominator of the Jacobi test by tables is at most $\mathrm{mulModLimit} = 2^{32} - 1$,
    where \verb|RP_MULDIV| is exact; the valuations and the slopes of the records of the divisors,
    \verb|int|s in the C code, as natural numbers, which they fit for coefficients below
    $2^{2^{31} - 1}$;
  \item the denominators of a shape listed, then tested and sifted one after the other; no blocks,
    and no threads, which deliver the blocks in their order; no correction of the moduli between
    the blocks (stage~I); no search intervals and no range of the denominators other than $1$ to $H$
    (stage~H);
  \item no failure for lack of memory, where \verb|sieving_info| drops the primes beyond the table
    whose words it cannot store; no output under \verb|DEBUG| or \verb|RATPOINTS_VERBOSE|.
\end{itemize}
The lines of this part of the C code that are neither modelled nor left to a later stage are
excluded in \verb|correspondence/ratpoints.toml|: the output under \verb|DEBUG| that stands between
statements that are not modelled, the management of memory (the arrays of the run and the worker
with the failures of their allocation, the release of the run and of the worker, the enlargement of
the buffer of the tables when \verb|sieving_info| has looked at more primes), and the pool of
threads with the threaded run; the output under \verb|RATPOINTS_VERBOSE| and \verb|DEBUG| inside a
modelled block goes with the block, whose deviation says so. Left to later stages are the packing
of the classes with their row offsets (stage~F3), the search intervals of the run (stage~H), and the
blocks with the correction of the moduli between them and the choice of the number of primes
examined (stage~I).

\section{Files and names}

The constants of the stage are in \verb|Ratpoints/Model/Constants.lean|, each citing its macro, and
no proof depends on their values: $\mathrm{maxDivisors}$, $\mathrm{veryBig}$,
$\mathrm{jacobiPrimes}$, $\mathrm{jacobiTable}$ and $\mathrm{mulModLimit}$, of which the proofs use
only $\mathrm{mulModLimit} < 2^{64}$ and $\mathrm{mulModLimit} < 2^{32}$
(\verb|Ratpoints.mulModLimit_lt_size| and \verb|Ratpoints.mulModLimit_lt_two_pow_32|, next to the
constant), so that a denominator of the Jacobi test by tables is a word, and a value at most
$\mathrm{mulModLimit}$ times a modulus $p \le 2^{32}$ is below $2^{64}$, where the reductions by a
reciprocal are exact (Lemma~\ref{lem:mulMod}). The valuation and the divisors of the leading
coefficient are in \verb|Ratpoints/Model/Setup.lean|, in the namespace \verb|Ratpoints.Setup|; the
2-adic information, the forbidden divisors, the Jacobi test, the shapes and the run in
\verb|Ratpoints/Model/Denominators.lean|, in the namespace \verb|Ratpoints.Sieve|; the class word in
the components of \verb|Ratpoints/Model/Sieve.lean| that make a piece and the numerators of one
denominator; and the steps of the search in \verb|Ratpoints/Model/FindPoints.lean|. The contracts
and the lemmas are in the files of \verb|Ratpoints/Proof| that the stage adds: the valuation in
\verb|Valuation.lean|; the 2-adic information in \verb|TwoAdic.lean|; the denominators of odd
degree, the factorization and the divisors with their slopes in \verb|Divisors.lean|; the forbidden
valuations and divisors in \verb|Forbidden.lean|; the Jacobi test by the binary algorithm in
\verb|Jacobi.lean| and by tables in \verb|JacobiTables.lean|; the lists and the tests of the shapes
in \verb|Shapes.lean|; and the run over a shape in \verb|Denominators.lean|. Unlike those of stages
C1 to D, the contracts of this stage are proved in files with Mathlib, next to the lemmas they
serve, since several of them speak of notions of Mathlib (the primes and their valuations, the
squares modulo a prime, the Jacobi symbol). The class word is in the contracts of
\verb|Ratpoints/Proof/Words.lean| and, since stage~F2, of \verb|Ratpoints/Proof/Sift.lean|, which
holds those from the walk on; \verb|Ratpoints/Proof/Words.lean| also holds the invariant of a loop
over the set bits of a word that the walk of the check, the 2-adic information and the words of the
shapes share; the filter of the reference by a test on the numerators is in
\verb|Ratpoints/Proof/Sieve.lean|, the denominators of a point and the filter of the reference by
the denominators in \verb|Ratpoints/Proof/BruteForce.lean|, the facts about the table of primes that
the factorization uses, checked by \verb|decide|, in \verb|Ratpoints/Proof/Params.lean|, the lemmas
on words, lists and ranges in core's namespaces in \verb|Ratpoints/Proof/Bits.lean|, the binary form
at $a = k b$, the squares modulo $p$ and the negation of the quarter of \verb|LONG_MIN| in
\verb|Ratpoints/Proof/Tables.lean|, and the denominators of the search with the theorems in
\verb|Ratpoints/Proof/FindPoints.lean|. The lemmas of Mathlib's generality are in
\verb|Ratpoints/Mathlib|, in the files that mirror their homes in Mathlib:
\verb|Algebra/BigOperators/Group/List/Basic|, \verb|Algebra/Group/Even|, \verb|Data/Int/GCD|,
\verb|Data/List/Basic|, \verb|Data/List/Pairwise|, \verb|Data/Nat/Factors| (with the part of a
number prime to another, \verb|Nat.coprimePart|, which is $b^*$), \verb|Data/Nat/PadicValNat|,
\verb|Data/Nat/Prime/Defs|, \verb|Data/Nat/Squarefree|, \verb|Data/ZMod/Basic|,
\verb|NumberTheory/LegendreSymbol/JacobiSymbol| and \verb|NumberTheory/Padics/PadicVal/Basic|; that
the macro \verb|CEIL| rounds up, a fact about the C code, is a private lemma next to its use.

\section{The evidence}\label{sec:den-evidence}

Besides the proofs:
\begin{itemize}
  \item a reading of the C code of this chapter at the pin, each part with a second reader who
    checked the reading: the rule of the 2-adic words, checked on 6367 curves over all 4096 pairs of
    residues, with the classes packed by strides of at most $8$; a build with the stride~$1$ for
    every class, which prints what the release prints on 2500 runs; the order of the output
    predicted byte for byte on 10088 runs (4836 curves; 70 refusals, 2085 runs whose denominators do
    not ascend), with 180 runs at heights up to 100000, 397 with a range of the denominators and 660
    with blocks of one step and three or four threads, all agreeing; the tests on the denominators
    of a transcription against the C code, test by test, on 956100 denominators of 3000 curves and
    2.64 million more at heights up to 150000; their soundness against the brute force on 2863
    curves of even degree with 25329 points and on 2993 curves with points of the valuations $1$ to
    $7$; and the three implementations of the Jacobi test agreeing on 31.6 million evaluations; no
    defect of the C code was found;
  \item before the stage, a sketch of the model in Lean, checked against an instrumented build of
    the C code, line by line, on 6260 runs and against the C program on 2540 runs of its own, and a
    prototype of the proofs for it, which found that the bound $m < 63$ of the C code on the masks
    of the forbidden valuations is not sound beyond the \verb|long|s (for $c = (3, 0, 3^{128})$ and
    $H \ge 3^{64}$ the prime $3$ would be tested with words, although $(1, 2 \cdot 3^{64}, 3^{64})$
    is a point), so that the masks are natural numbers; a review of the C ranges of the sketch (34
    findings, all confirmed) and of the plan of this chapter by four readers (53 findings, 39
    confirmed and 12 in part, all applied), each finding checked by a second reader; after the
    stage, two independent reviews, each finding checked by a second reader: of the C ranges, the
    deviations and the exclusions against the C code (9 findings, all confirmed: the definitions of
    \verb|RP_MULDIV| go to the Jacobi test by tables, and eight deviations name the types and
    arguments they left out), and of the code against the conventions of the project (85 findings,
    82 confirmed), applied in eighteen rounds, among them one invariant for the loops over the set
    bits of a word, a predicate for the result of the subtraction loop of the Jacobi symbol, the
    part of a number prime to another as a definition, unfolding lemmas for the tests, the lemmas of
    Mathlib's generality in their homes and at their generality, the weakest hypotheses, and names
    from the statements;
  \item the \verb|#guard| examples of the model: the valuation, the factorizations with their
    cofactors, the divisors and the slopes against what \verb|ratpoints -v| and the build with
    \verb|DEBUG| print, the boundary of $512$ divisors, the words of the 2-adic information, the
    forbidden divisors, the tables of the Jacobi test and the denominators sifted against the
    instrumented build, the binary algorithm against Euler's criterion and the test by tables
    against the binary algorithm, and the search with the divisor shape and at the 2-adic exit
    against ratpoints;
  \item the differential tests of Section~\ref{sec:denCli}: 3219 comparisons in the default suites
    and 13 in \verb|test4high|, none failed, every point list compared in its order;
  \item Comparator accepted the solution of the challenge, with the model of this stage, on
    2026-10-12;
  \item the project builds, with all its proofs, for other values of the constants of the build, the
    five new ones among them; so no proof depends on their values;
  \item every proof of the stage takes at most 20 lines.
\end{itemize}
